\chapter*{Avant-propos}
\addcontentsline{toc}{chapter}{Avant-propos}

Ce manuscrit décrit \reactant{}, un analyseur statique de programmes React
fondé sur l'interprétation abstraite, tel qu'il existe au commit
\snapshotCommit{} du dépôt (\snapshotDate). Il a été écrit pour qu'un lecteur
qui connaît Rust et React côté utilisateur, sans formation préalable en analyse
statique, en sorte en connaissant très bien la codebase : ses fichiers, ses
types, ses algorithmes, ses décisions d'architecture et ses limites.

\medskip

\noindent\textbf{Parcours.} L'ordre des chapitres va du concret vers
l'abstrait. La première partie pose le problème : les bugs React visés, le
modèle d'exécution de React que l'on abstrait, les rappels d'interprétation
abstraite nécessaires, puis un tour d'architecture de la codebase. Les parties
suivantes suivent un fichier source dans l'ordre où l'analyseur le traite : le
parseur et la détection des composants, la représentation intermédiaire et le
lowering, les domaines abstraits, le moteur de point fixe et ses relations,
l'analyse inter-composants, puis les règles qui transforment ces faits en
diagnostics. La dernière partie décrit l'outil lui-même (ligne de commande,
distribution, méthode de mesure), l'histoire de ses décisions et ses limites
connues. Trois annexes servent de référence : les fiches des décisions
d'architecture, le glossaire et le catalogue des règles. La carte de la
codebase prévue en quatrième annexe n'a pas été rédigée dans cette version ; le
tour d'architecture (chapitre~\ref{chap:tour-architecture}) en tient lieu.

Un lecteur pressé peut lire les chapitres 1 à 4 puis piocher ; un lecteur qui
veut contribuer au moteur lira dans l'ordre jusqu'au chapitre 16 ; un lecteur
qui veut écrire une règle lira les chapitres 1 à 4, 17 et 22.

\medskip

\noindent\textbf{Conventions.} Les extraits de code Rust sont pris tels quels
dans le dépôt, avec leurs numéros de ligne réels ; les programmes React
d'exemple sont donnés en TSX et les sorties de l'analyseur ont toutes été
observées sur le binaire construit au commit photographié. Les encadrés
\emph{Rappel React} rappellent un point de la sémantique de React nécessaire à
la compréhension ; les encadrés \emph{Décision} résument une décision
d'architecture (ADR) et les alternatives refusées ; les encadrés \emph{Piège}
signalent une subtilité ou une erreur classique ; les encadrés
\emph{Exercice} proposent un travail au lecteur, avec solution. Les
identifiants, noms de fichiers et noms de règles restent en anglais, tels qu'ils
apparaissent dans le code ; les notions mathématiques sont en français avec une
notation unifiée (ordre $\lleq$, join $\join$, meet $\meet$, widening $\widen$,
concrétisation $\gam$).

\medskip

\noindent\textbf{État du texte.} Chaque chapitre a été rédigé à partir de
dossiers techniques établis sur le code, puis relu par un vérificateur qui a
recontrôlé chaque extrait et rejoué chaque sortie. Les constats faits en cours
de rédaction sur le code lui-même (faux négatifs ou faux positifs observés,
écarts entre une ADR et le code) sont rassemblés au
chapitre~\ref{chap:limites-perspectives} et présentés comme des observations
au commit \snapshotCommit{}, non comme des jugements. Une passe finale
d'harmonisation reste à faire ; les renvois croisés et l'index peuvent encore
comporter des irrégularités.

\medskip

\noindent\textbf{Sources.} La documentation officielle de React
\cite{ReactDocs}, la sémantique formelle React-tRace \cite{LeeAhnYi2025} et
les travaux fondateurs de l'interprétation abstraite
\cite{CousotCousot77,CousotCousot79,RivalYi2020} sont les références externes
de ce manuscrit ; tout le reste vient du dépôt \cite{ReactantRepo}.
